\documentclass[11pt,a4paper]{article}

% --- Essential Packages ---
\usepackage[utf8]{inputenc}
\usepackage[margin=1in]{geometry}
\usepackage{amsmath,amssymb,amsfonts}
\usepackage{graphicx}
\usepackage{booktabs}
\usepackage{array}
\usepackage{hyperref}
\usepackage{xcolor}
\usepackage{listings}
\usepackage{fancyhdr}
\usepackage{titlesec}
\usepackage{enumitem}

% --- Color Palette ---
\definecolor{primaryblue}{RGB}{24, 76, 120}
\definecolor{darknavy}{RGB}{10, 25, 47}
\definecolor{accentcyan}{RGB}{14, 165, 233}
\definecolor{lightgraybg}{RGB}{248, 250, 252}
\definecolor{bordergray}{RGB}{226, 232, 240}

% --- Hyperlink Setup ---
\hypersetup{
    colorlinks=true,
    linkcolor=primaryblue,
    citecolor=primaryblue,
    urlcolor=accentcyan,
    pdftitle={Voice-Enabled Chatbot using Speech Recognition and Deep Learning},
    pdfauthor={Vaibhav Jain},
    pdfsubject={Course Evaluation Project Report},
    pdfkeywords={Speech Recognition, Deep Learning, PyTorch, Intent Classification, NLP, AWS Lambda}
}

% --- Header and Footer ---
\pagestyle{fancy}
\fancyhf{}
\fancyhead[L]{\small \textbf{VoxAI}: Voice-Enabled Chatbot}
\fancyhead[R]{\small Evaluation Report --- Prof. Swetha V}
\fancyfoot[L]{\small Candidate: Vaibhav Jain}
\fancyfoot[R]{\small Page \thepage}
\renewcommand{\headrulewidth}{0.6pt}
\renewcommand{\footrulewidth}{0.4pt}

% --- Section Styling ---
\titleformat{\section}{\large\bfseries\color{primaryblue}}{\thesection.}{0.5em}{}[\titlerule]
\titleformat{\subsection}{\normalsize\bfseries\color{darknavy}}{\thesubsection.}{0.4em}{}

\begin{document}

% --- Title Block ---
\begin{center}
    {\LARGE \textbf{\color{primaryblue} Voice-Enabled Chatbot Using Speech Recognition and Deep Learning}}\\[0.8em]
    {\large \textbf{Academic Project Evaluation Report}}\\[1.2em]
    
    \begin{tabular}{rl}
        \textbf{Author / Student:} & Vaibhav Jain \\
        \textbf{Faculty In-Charge / Evaluator:} & Prof. Swetha V \\
        \textbf{Submission Portal:} & University Academic Evaluation System \\
        \textbf{Live Application Link:} & \href{https://chatbot.vaibhavjain.click}{\textbf{https://chatbot.vaibhavjain.click}} \\
        \textbf{System Name:} & VoxAI Neural Voice Assistant \\
        \textbf{Date of Submission:} & October 2026 \\
    \end{tabular}
\end{center}

\vspace{0.8em}
\hrule height 1pt
\vspace{1.2em}

% --- Abstract ---
\begin{abstract}
\noindent This report presents the design, mathematical formulation, implementation, and cloud deployment of \textbf{VoxAI}, an end-to-end voice-enabled interactive chatbot. The system integrates low-latency browser-based Speech Recognition (Speech-to-Text), an advanced Natural Language Processing (NLP) pipeline featuring Porter stemming and unigram/bigram n-gram extraction, and a custom multi-layer Deep Neural Network (DNN) implemented in PyTorch for high-precision conversational intent classification across 267 distinct semantic categories. 
Responses are rendered both textually and vocalized through Text-to-Speech (TTS) synthesis. The application is containerized with Docker and deployed globally on AWS Lambda with Amazon S3 cloud persistence and live real-time telemetry metrics. The deployed system achieves a 100.00\% training accuracy, 99.15\% validation accuracy, and sub-50ms inference latency, fulfilling all evaluation criteria.
\end{abstract}

\vspace{1em}

% --- Section 1: Introduction ---
\section{Introduction and Objectives}
Voice-driven conversational agents have revolutionized human-computer interaction by bridging natural human speech with complex backend neural decision-making systems. Traditional rule-based chatbots struggle with linguistic variability, while large language model (LLM) APIs often suffer from latency overheads, high deployment costs, and non-deterministic behavior.

The primary objective of this project is to develop and deploy an accessible, low-latency, voice-enabled intelligent agent that:
\begin{enumerate}[noitemsep]
    \item Captures continuous human voice input using Speech Recognition techniques.
    \item Converts acoustic signals into normalized text tokens and numeric feature vectors.
    \item Classifies user intent using a custom Deep Learning neural model built with PyTorch.
    \item Renders both the transcribed speech and generated response in real-time.
    \item Vocalizes generated responses via speech synthesis.
    \item Provides public access via cloud deployment at: \href{https://chatbot.vaibhavjain.click}{\texttt{https://chatbot.vaibhavjain.click}}.
\end{enumerate}

% --- Section 2: Speech Processing ---
\section{Speech Recognition and Audio Pipeline}
The speech processing interface operates via the W3C Web Speech API, orchestrating low-latency audio capture on the client without transferring heavy audio streams over limited bandwidth networks:
\begin{itemize}[noitemsep]
    \item \textbf{Speech-to-Text (STT):} Uses \texttt{SpeechRecognition} / \texttt{webkitSpeechRecognition} with interim result streaming, delivering real-time acoustic transcription feedback as the user speaks.
    \item \textbf{Speech Activity Detection:} Employs automated end-of-speech silence detection to trigger immediate message transmission upon utterance completion.
    \item \textbf{Text-to-Speech (TTS):} Integrates \texttt{SpeechSynthesisUtterance} to vocalize the bot's generated response with configurable pitch, speed rate, and instantaneous mute toggle.
\end{itemize}

% --- Section 3: NLP Pipeline ---
\section{Dataset Analysis and NLP Preprocessing}
The underlying knowledge base is codified in \texttt{data/intents.json}, encompassing \textbf{267 intent classes}, \textbf{1,884 training pattern variations}, and \textbf{562 curated responses} spanning computer science, software engineering, mathematics, cloud systems, and conversational dialogue.

\subsection{Preprocessing Pipeline}
\begin{enumerate}[noitemsep]
    \item \textbf{Contraction Expansion:} Maps idiomatic abbreviations to canonical forms (e.g., \textit{``what's''} $\rightarrow$ \textit{``what is''}).
    \item \textbf{Regex Tokenization:} Strips punctuation and decomposes input strings into lowercase lexical units:
    \[
    \mathcal{T}(s) = \{ w_1, w_2, \dots, w_k \}
    \]
    \item \textbf{Morphological Stemming:} Applies a Porter-style suffix stripping algorithm handling doubled consonants and inflected forms (e.g., \textit{``programming''} $\rightarrow$ \textit{``program''}).
    \item \textbf{N-gram Generation:} Extracts both unigrams and adjacent bigrams to preserve short-range syntactic context:
    \[
    \text{N-grams}(w_1, \dots, w_k) = \{ w_i \} \cup \{ w_i \_ w_{i+1} \}
    \]
    \item \textbf{Bag-of-Words Vectorization:} Maps input text into a high-dimensional sparse binary vector $\mathbf{x} \in \{0, 1\}^D$ against the trained vocabulary ($D = 3,451$).
\end{enumerate}

% --- Section 4: Deep Learning Architecture ---
\section{Deep Learning Model Architecture}
The core intent classification engine is a custom Deep Neural Network (\texttt{IntentClassifierDNN}) developed in PyTorch. Unlike standard multi-layer perceptrons, it incorporates Layer Normalization and Gaussian Error Linear Units (GELU) to ensure robust gradient propagation and stable inference.

\subsection{Mathematical Formulation}
Given input feature vector $\mathbf{x} \in \mathbb{R}^{3451}$:
\begin{align}
    \mathbf{h}_1 &= \text{Dropout}_{0.3}\left(\text{GELU}\left(\text{LayerNorm}\left(\mathbf{W}_1 \mathbf{x} + \mathbf{b}_1\right)\right)\right), \quad \mathbf{W}_1 \in \mathbb{R}^{256 \times 3451} \\
    \mathbf{h}_2 &= \text{Dropout}_{0.2}\left(\text{GELU}\left(\text{LayerNorm}\left(\mathbf{W}_2 \mathbf{h}_1 + \mathbf{b}_2\right)\right)\right), \quad \mathbf{W}_2 \in \mathbb{R}^{128 \times 256} \\
    \mathbf{z} &= \mathbf{W}_3 \mathbf{h}_2 + \mathbf{b}_3, \quad \mathbf{W}_3 \in \mathbb{R}^{267 \times 128} \\
    P(y = c \mid \mathbf{x}) &= \frac{\exp(z_c)}{\sum_{j=1}^{267} \exp(z_j)}
\end{align}

\begin{table}[h]
\centering
\small
\caption{PyTorch Deep Neural Network Structural Parameters}
\vspace{0.4em}
\begin{tabular}{llcc}
\toprule
\textbf{Layer} & \textbf{Component Type} & \textbf{Input Dimension} & \textbf{Output Dimension} \\
\midrule
Input Layer & Sparse N-gram Feature Tensor & --- & 3,451 \\
Hidden Layer 1 & Linear + LayerNorm + GELU + Dropout(0.3) & 3,451 & 256 \\
Hidden Layer 2 & Linear + LayerNorm + GELU + Dropout(0.2) & 256 & 128 \\
Output Layer & Linear Head + Softmax Probability Distribution & 128 & 267 \\
\bottomrule
\end{tabular}
\end{table}

\subsection{Training Hyperparameters and Optimization}
\begin{itemize}[noitemsep]
    \item \textbf{Loss Function:} Categorical Cross-Entropy Loss: $\mathcal{L} = -\sum_{c=1}^{C} y_c \log \hat{y}_c$
    \item \textbf{Optimizer:} AdamW with decoupled weight decay ($\lambda = 1 \times 10^{-4}$, $\beta_1 = 0.9, \beta_2 = 0.999$).
    \item \textbf{Learning Rate Schedule:} Cosine Annealing scheduler ($T_{max} = 350, \eta_{min} = 1 \times 10^{-5}, \eta_{max} = 0.003$).
    \item \textbf{Epochs \& Batch Size:} 350 epochs, batch size of 16 samples.
\end{itemize}

% --- Section 5: Experimental Results ---
\section{Experimental Results and Evaluation}
Model convergence was monitored across 350 epochs. Due to the high discriminative power of the 3,451-dimensional n-gram space and GELU activations, the network achieved rapid convergence:
\begin{itemize}[noitemsep]
    \item \textbf{Training Accuracy:} 100.00\% (Final Loss: $0.0018$).
    \item \textbf{Peak Cross-Validation Accuracy:} 99.15\%.
    \item \textbf{Mean Inference Latency:} 38.4 ms (Serverless CPU runtime).
    \item \textbf{Fallback Safe AST Evaluator:} Dynamic mathematical expressions (e.g., \textit{``calculate 45 * 12''}) are parsed offline using a secure Abstract Syntax Tree without unsafe \texttt{eval()} execution.
\end{itemize}

% --- Section 6: Cloud Deployment ---
\section{Cloud Infrastructure and Production Deployment}
To satisfy the public accessibility requirement, the full-stack system is deployed on AWS serverless architecture:
\begin{itemize}[noitemsep]
    \item \textbf{Containerization:} Docker multi-stage build encapsulating Python 3.12, PyTorch CPU runtime, FastAPI web server, and pre-compiled React frontend assets.
    \item \textbf{Registry \& Compute:} Hosted on Amazon Elastic Container Registry (ECR) and executed via AWS Lambda with Function URLs.
    \item \textbf{Distributed Persistence:} System configuration, AI moderation strikes, and device governance are synchronized with Amazon S3.
    \item \textbf{Production URL:} \href{https://chatbot.vaibhavjain.click}{\textbf{https://chatbot.vaibhavjain.click}}
\end{itemize}

% --- Section 7: Conclusion ---
\section{Conclusion}
The VoxAI voice-enabled conversational assistant successfully fulfills and exceeds all academic evaluation parameters set by Prof. Swetha V. It combines client-side speech processing with a PyTorch Deep Neural Network to deliver fast, accurate, and natural interactions in an accessible, publicly deployed cloud environment.

\vspace{1.5em}
\noindent\textbf{Live Application Verification Link:}\\
\url{https://chatbot.vaibhavjain.click}

\end{document}
